\section{Conclusion}
\label{sec:conclusion}

The engineering of the CropSense platform marks a significant transition from theoretical computer vision experiments to a fully operational agricultural application. Building a robust system capable of handling the extreme variations in field topologies, aggressive lighting gradients, and densely packed clusters of the GWHD demanded a strictly deterministic approach.

Our initial evaluations proved that off-the-shelf object detectors are too fragile for real-world agronomy. They suffer from scale annihilation when attempting to detect distant crops and exhibit a severe tendency to hallucinate false positives when confused by ground shadows or overlapping leaves. However, by substituting blind parameter sweeps with a highly focused, error-driven optimization cycle, we were able to completely rewire the network's behavior. We traded a negligible drop in absolute recall for a massive +39.4\% surge in prediction certainty, converting the model from a reckless guesser into a highly decisive, operationally stable engine.

Despite achieving these milestones, our forensic analysis exposed a clear biological constraint. The current iteration of CropSense fails consistently when presented with entirely green, immature wheat canopies because the core dataset is heavily biased toward mature, yellowing phenotypes. 

Therefore, future enhancement of the system must prioritize expanding the training corpus to cover the full spectrum of the BBCH phenological growth scale. Additionally, incorporating multi-angle UAV imagery rather than relying exclusively on top-down nadir angles will improve the 3D spatial robustness of the model. Ultimately, by seamlessly integrating high-confidence computer vision with automated data reporting, CropSense demonstrates the immense, scalable potential of AI in safeguarding modern agricultural yields.